\documentclass[11pt]{article}
\usepackage[margin=1in]{geometry}
\usepackage{amsmath, amssymb, amsthm}
\usepackage{algorithm}
\usepackage{algpseudocode}
\usepackage{bm}

\newtheorem{theorem}{Theorem}
\newtheorem{lemma}{Lemma}
\newtheorem{assumption}{Assumption}

\title{Comprehensive Convergence Analysis of RMSProp \\ for Non-Convex Smooth Functions}
\author{}
\date{}

\begin{document}

\maketitle

\section*{Algorithm Name}
\textbf{RMSProp} (Root Mean Square Propagation)

\section*{Mathematical Formulation}
RMSProp is an adaptive learning rate optimization algorithm designed to resolve the radically diminishing learning rates of AdaGrad. It achieves this by utilizing a moving average of the squared gradients to normalize the gradient updates. 

Let $f: \mathbb{R}^d \rightarrow \mathbb{R}$ be the objective function to be minimized. At each iteration $t$, let $g_t$ be the stochastic gradient of $f$ evaluated at the current parameters $w_t$. The algorithm updates a moving average of the squared gradients, $v_t \in \mathbb{R}^d$, and uses it to scale the step size element-wise.

The formal update rules are defined element-wise as follows:
\begin{align*}
    g_t &= \nabla f(w_t; \xi_t) \\
    v_t &= \beta v_{t-1} + (1 - \beta) g_t^2 \\
    w_{t+1} &= w_t - \frac{\alpha}{\sqrt{v_t} + \epsilon} \odot g_t
\end{align*}
where:
\begin{itemize}
    \item $\alpha > 0$ is the base learning rate (step size).
    \item $\beta \in [0, 1)$ is the exponential decay rate for the second moment estimate.
    \item $\epsilon > 0$ is a small scalar added for numerical stability and strict theoretical bounds.
    \item $\odot$ denotes element-wise multiplication, and the square/square-root operations on vectors are also applied element-wise.
\end{itemize}

\section*{Pseudocode}
\begin{algorithm}[H]
\caption{RMSProp Optimization Algorithm}
\begin{algorithmic}[1]
\Require Base step size $\alpha > 0$, decay rate $\beta \in [0, 1)$, stability parameter $\epsilon > 0$
\State Initialize parameters $w_1 \in \mathbb{R}^d$
\State Initialize moving average $v_0 = \mathbf{0} \in \mathbb{R}^d$
\For{$t = 1, 2, \dots, T$}
    \State Sample data $\xi_t$ and compute stochastic gradient $g_t = \nabla f(w_t; \xi_t)$
    \State Update second moment estimate: $v_t = \beta v_{t-1} + (1 - \beta) g_t \odot g_t$
    \State Compute adaptive learning rate: $\eta_t = \frac{\alpha}{\sqrt{v_t} + \epsilon}$ \Comment{Element-wise operations}
    \State Update parameters: $w_{t+1} = w_t - \eta_t \odot g_t$
\EndFor
\State \textbf{return} $w_{T+1}$ (or a randomly chosen $w_t$ from the trajectory)
\end{algorithmic}
\end{algorithm}

\section*{Dry Run Example}
Consider minimizing the scalar quadratic function $f(w) = (w - 3)^2$. The exact gradient is $g_t = \nabla f(w_t) = 2(w_t - 3)$. 
We initialize $w_1 = 0$, $v_0 = 0$. Let the hyperparameters be $\alpha = 0.1$, $\beta = 0.9$, and $\epsilon = 10^{-8}$.

\textbf{Iteration 1 ($t=1$):}
\begin{itemize}
    \item Current parameter: $w_1 = 0$
    \item Gradient: $g_1 = 2(0 - 3) = -6$
    \item Second moment: $v_1 = 0.9(0) + 0.1(-6)^2 = 3.6$
    \item Update: $w_2 = w_1 - \frac{0.1}{\sqrt{3.6} + 10^{-8}} (-6) \approx 0 - \frac{0.1}{1.897366} (-6) \approx 0.3162$
\end{itemize}

\textbf{Iteration 2 ($t=2$):}
\begin{itemize}
    \item Current parameter: $w_2 = 0.3162$
    \item Gradient: $g_2 = 2(0.3162 - 3) = -5.3676$
    \item Second moment: $v_2 = 0.9(3.6) + 0.1(-5.3676)^2 = 3.24 + 2.8811 = 6.1211$
    \item Update: $w_3 = 0.3162 - \frac{0.1}{\sqrt{6.1211} + 10^{-8}} (-5.3676) \approx 0.3162 + \frac{0.53676}{2.4741} \approx 0.5332$
\end{itemize}

\textbf{Iteration 3 ($t=3$):}
\begin{itemize}
    \item Current parameter: $w_3 = 0.5332$
    \item Gradient: $g_3 = 2(0.5332 - 3) = -4.9336$
    \item Second moment: $v_3 = 0.9(6.1211) + 0.1(-4.9336)^2 = 5.5090 + 2.4340 = 7.9430$
    \item Update: $w_4 = 0.5332 - \frac{0.1}{\sqrt{7.9430} + 10^{-8}} (-4.9336) \approx 0.5332 + \frac{0.49336}{2.8183} \approx 0.7083$
\end{itemize}
Notice how the parameter $w$ steadily moves towards the optimum $w^* = 3$, while the effective step size adapts based on the accumulated gradient magnitudes.

\section*{Assumptions}
To rigorously prove the convergence of RMSProp for non-convex functions, we require the following standard assumptions:

\begin{assumption}[$L$-Smoothness]
The objective function $f$ is differentiable and $L$-smooth, meaning its gradients are Lipschitz continuous:
\begin{equation*}
    \|\nabla f(x) - \nabla f(y)\| \leq L \|x - y\|, \quad \forall x, y \in \mathbb{R}^d
\end{equation*}
This implies the standard quadratic upper bound: 
$f(y) \leq f(x) + \langle \nabla f(x), y - x \rangle + \frac{L}{2} \|y - x\|^2$.
\end{assumption}

\begin{assumption}[Lower Boundedness]
The function $f$ is bounded below by some finite value $f^*$:
\begin{equation*}
    f(w) \geq f^*, \quad \forall w \in \mathbb{R}^d
\end{equation*}
\end{assumption}

\begin{assumption}[Unbiased Stochastic Gradients]
The stochastic gradient $g_t$ is an unbiased estimator of the true gradient:
\begin{equation*}
    \mathbb{E}[g_t \mid w_t] = \nabla f(w_t)
\end{equation*}
\end{assumption}

\begin{assumption}[Bounded Variance]
The variance of the stochastic gradient is bounded:
\begin{equation*}
    \mathbb{E}[\|g_t - \nabla f(w_t)\|^2 \mid w_t] \leq \sigma^2
\end{equation*}
\end{assumption}

\begin{assumption}[Bounded Gradients]
The stochastic gradients are bounded almost surely in the infinity norm by a constant $G$:
\begin{equation*}
    \|g_t\|_\infty \leq G, \quad \text{almost surely.}
\end{equation*}
(This is a standard assumption in the analysis of adaptive methods like Adam and RMSProp to ensure the preconditioner $v_t$ does not grow unboundedly).
\end{assumption}

\section*{Main Convergence Theorem}
\begin{theorem}
Under Assumptions 1-5, if the learning rate $\alpha$ is chosen such that $\alpha \leq \frac{\epsilon}{L}$, then the sequence of iterates $\{w_t\}_{t=1}^T$ generated by RMSProp satisfies:
\begin{equation*}
    \frac{1}{T} \sum_{t=1}^T \mathbb{E}[\|\nabla f(w_t)\|^2] \leq \frac{2(G + \epsilon)(f(w_1) - f^*)}{\alpha T} + \frac{(G + \epsilon)\sigma^2}{\epsilon}
\end{equation*}
Furthermore, by setting $\alpha = \min\left(\frac{\epsilon}{L}, \frac{C}{\sqrt{T}}\right)$ for some constant $C > 0$, the algorithm achieves a convergence rate of $\mathcal{O}\left(\frac{1}{\sqrt{T}}\right)$ to a stationary point.
\end{theorem}

\section*{Theoretical Properties}
\begin{enumerate}
    \item \textbf{Adaptive Scaling:} RMSProp normalizes the gradient by a moving average of its recent magnitude. This ensures that parameters with sparse gradients take larger steps, while parameters with dense, large gradients take smaller, controlled steps.
    \item \textbf{Bounded Preconditioner:} Due to the convex combination in the update rule $v_t = \beta v_{t-1} + (1-\beta)g_t^2$, if gradients are bounded by $G$, the components of $v_t$ are strictly bounded by $G^2$.
    \item \textbf{Non-vanishing Step Sizes:} Unlike AdaGrad where $v_t$ grows monotonically (causing premature convergence), RMSProp's exponential decay ensures $v_t$ stays bounded, allowing the algorithm to continue learning indefinitely.
\end{enumerate}

\section*{Discussion}
The inclusion of $\epsilon$ is not merely a practical hack to avoid division by zero; it strictly controls the maximum possible effective step size. In our proof, $\epsilon$ dictates the upper bound on the variance introduced by the adaptive denominator. If $\alpha$ is kept constant, the algorithm converges to a noise ball determined by $\frac{(G+\epsilon)\sigma^2}{\epsilon}$. To achieve exact convergence to a stationary point (zero gradient), the learning rate $\alpha$ must be annealed at a rate of $\mathcal{O}(1/\sqrt{T})$.

\section*{Preliminaries (Definitions and Lemmas)}

\begin{lemma}[Algebraic Identity]
For any two real numbers $a$ and $b$, the following identity holds:
\begin{equation*}
    -a b = \frac{1}{2} \left( (a - b)^2 - a^2 - b^2 \right)
\end{equation*}
\end{lemma}
\begin{proof}
Expand the right-hand side: $\frac{1}{2}(a^2 - 2ab + b^2 - a^2 - b^2) = \frac{1}{2}(-2ab) = -ab$.
\end{proof}

\begin{lemma}[Boundedness of $v_t$]
Under Assumption 5 ($\|g_t\|_\infty \leq G$) and initializing $v_0 = 0$, for all $t \geq 1$ and all coordinates $i \in \{1, \dots, d\}$, we have $v_{t,i} \leq G^2$.
\end{lemma}
\begin{proof}
We proceed by induction. Base case $v_0 = 0 \leq G^2$. Assume $v_{t-1,i} \leq G^2$. Then:
$v_{t,i} = \beta v_{t-1,i} + (1-\beta)g_{t,i}^2 \leq \beta G^2 + (1-\beta)G^2 = G^2$.
\end{proof}

\section*{Main Proof (Step-by-Step Derivation)}

Let $g_{t,i}$ and $\nabla f(w_t)_i$ denote the $i$-th coordinate of the stochastic gradient and the true gradient, respectively.

\begin{align}
    \text{By the $L$-smoothness of } f, \text{ we have:} \nonumber \\
    f(w_{t+1}) &\leq f(w_t) + \langle \nabla f(w_t), w_{t+1} - w_t \rangle + \frac{L}{2} \|w_{t+1} - w_t\|^2 \tag{Step 1} \\
    \text{Substitute the RMSProp update } w_{t+1} - w_t = - \alpha \frac{g_t}{\sqrt{v_t} + \epsilon}: \nonumber \\
    f(w_{t+1}) &\leq f(w_t) - \alpha \sum_{i=1}^d \frac{\nabla f(w_t)_i g_{t,i}}{\sqrt{v_{t,i}} + \epsilon} + \frac{L \alpha^2}{2} \sum_{i=1}^d \frac{g_{t,i}^2}{(\sqrt{v_{t,i}} + \epsilon)^2} \tag{Step 2} \\
    \text{Apply Lemma 1 to the inner product term } a = \nabla f(w_t)_i, b = g_{t,i}: \nonumber \\
    - \frac{\nabla f(w_t)_i g_{t,i}}{\sqrt{v_{t,i}} + \epsilon} &= \frac{(\nabla f(w_t)_i - g_{t,i})^2}{2(\sqrt{v_{t,i}} + \epsilon)} - \frac{\nabla f(w_t)_i^2}{2(\sqrt{v_{t,i}} + \epsilon)} - \frac{g_{t,i}^2}{2(\sqrt{v_{t,i}} + \epsilon)} \tag{Step 3} \\
    \text{Substitute this identity back into the smoothness bound:} \nonumber \\
    f(w_{t+1}) &\leq f(w_t) + \frac{\alpha}{2} \sum_{i=1}^d \frac{(\nabla f(w_t)_i - g_{t,i})^2}{\sqrt{v_{t,i}} + \epsilon} - \frac{\alpha}{2} \sum_{i=1}^d \frac{\nabla f(w_t)_i^2}{\sqrt{v_{t,i}} + \epsilon} \nonumber \\
    &\quad - \frac{\alpha}{2} \sum_{i=1}^d \frac{g_{t,i}^2}{\sqrt{v_{t,i}} + \epsilon} + \frac{L \alpha^2}{2} \sum_{i=1}^d \frac{g_{t,i}^2}{(\sqrt{v_{t,i}} + \epsilon)^2} \tag{Step 4} \\
    \text{Group the terms involving } g_{t,i}^2: \nonumber \\
    - \frac{\alpha}{2} \frac{g_{t,i}^2}{\sqrt{v_{t,i}} + \epsilon} + \frac{L \alpha^2}{2} \frac{g_{t,i}^2}{(\sqrt{v_{t,i}} + \epsilon)^2} &= \frac{\alpha g_{t,i}^2}{2(\sqrt{v_{t,i}} + \epsilon)} \left( \frac{L \alpha}{\sqrt{v_{t,i}} + \epsilon} - 1 \right) \tag{Step 5} \\
    \text{Require } \alpha \leq \frac{\epsilon}{L}. \text{ Since } \sqrt{v_{t,i}} \geq 0 \implies \sqrt{v_{t,i}} + \epsilon \geq \epsilon, \text{ we have } \frac{L \alpha}{\sqrt{v_{t,i}} + \epsilon} \leq \frac{L \alpha}{\epsilon} \leq 1. \nonumber \\
    \text{Thus, the grouped term is non-positive and can be dropped:} \nonumber \\
    f(w_{t+1}) &\leq f(w_t) - \frac{\alpha}{2} \sum_{i=1}^d \frac{\nabla f(w_t)_i^2}{\sqrt{v_{t,i}} + \epsilon} + \frac{\alpha}{2} \sum_{i=1}^d \frac{(\nabla f(w_t)_i - g_{t,i})^2}{\sqrt{v_{t,i}} + \epsilon} \tag{Step 6} \\
    \text{By Lemma 2, } v_{t,i} \leq G^2, \text{ so } \sqrt{v_{t,i}} + \epsilon \leq G + \epsilon. \text{ This bounds the descent term:} \nonumber \\
    - \frac{\alpha}{2} \sum_{i=1}^d \frac{\nabla f(w_t)_i^2}{\sqrt{v_{t,i}} + \epsilon} &\leq - \frac{\alpha}{2(G + \epsilon)} \sum_{i=1}^d \nabla f(w_t)_i^2 = - \frac{\alpha}{2(G + \epsilon)} \|\nabla f(w_t)\|^2 \tag{Step 7} \\
    \text{Since } \sqrt{v_{t,i}} \geq 0, \text{ we have } \frac{1}{\sqrt{v_{t,i}} + \epsilon} \leq \frac{1}{\epsilon}. \text{ This bounds the variance term:} \nonumber \\
    \frac{\alpha}{2} \sum_{i=1}^d \frac{(\nabla f(w_t)_i - g_{t,i})^2}{\sqrt{v_{t,i}} + \epsilon} &\leq \frac{\alpha}{2\epsilon} \sum_{i=1}^d (\nabla f(w_t)_i - g_{t,i})^2 = \frac{\alpha}{2\epsilon} \|g_t - \nabla f(w_t)\|^2 \tag{Step 8} \\
    \text{Combine these bounds back into the main inequality:} \nonumber \\
    f(w_{t+1}) &\leq f(w_t) - \frac{\alpha}{2(G + \epsilon)} \|\nabla f(w_t)\|^2 + \frac{\alpha}{2\epsilon} \|g_t - \nabla f(w_t)\|^2 \tag{Step 9} \\
    \text{Take the unconditional expectation } \mathbb{E}[\cdot] \text{ and apply Assumption 4 } (\mathbb{E}[\|g_t - \nabla f\|^2] \leq \sigma^2): \nonumber \\
    \mathbb{E}[f(w_{t+1})] &\leq \mathbb{E}[f(w_t)] - \frac{\alpha}{2(G + \epsilon)} \mathbb{E}[\|\nabla f(w_t)\|^2] + \frac{\alpha \sigma^2}{2\epsilon} \tag{Step 10} \\
    \text{Rearrange to isolate the expected squared gradient norm:} \nonumber \\
    \frac{\alpha}{2(G + \epsilon)} \mathbb{E}[\|\nabla f(w_t)\|^2] &\leq \mathbb{E}[f(w_t)] - \mathbb{E}[f(w_{t+1})] + \frac{\alpha \sigma^2}{2\epsilon} \tag{Step 11} \\
    \text{Sum over } t = 1 \text{ to } T \text{ (the first two terms telescope):} \nonumber \\
    \frac{\alpha}{2(G + \epsilon)} \sum_{t=1}^T \mathbb{E}[\|\nabla f(w_t)\|^2] &\leq f(w_1) - \mathbb{E}[f(w_{T+1})] + \frac{T \alpha \sigma^2}{2\epsilon} \tag{Step 12} \\
    \text{Apply the lower bound } f(w_{T+1}) \geq f^* \text{ and divide by } T \frac{\alpha}{2(G + \epsilon)}: \nonumber \\
    \frac{1}{T} \sum_{t=1}^T \mathbb{E}[\|\nabla f(w_t)\|^2] &\leq \frac{2(G + \epsilon)(f(w_1) - f^*)}{\alpha T} + \frac{(G + \epsilon)\sigma^2}{\epsilon} \tag{Step 13}
\end{align}

\section*{Rate Derivation (With Constants)}
To establish the final convergence rate with respect to the horizon $T$, we must choose an appropriate learning rate $\alpha$. 
Let $\alpha = \min\left(\frac{\epsilon}{L}, \frac{C}{\sqrt{T}}\right)$ for some constant $C > 0$. 
For sufficiently large $T$, the condition $\alpha = \frac{C}{\sqrt{T}} \leq \frac{\epsilon}{L}$ holds. Substituting this into Step 13 yields:
\begin{align}
    \frac{1}{T} \sum_{t=1}^T \mathbb{E}[\|\nabla f(w_t)\|^2] &\leq \frac{2(G + \epsilon)(f(w_1) - f^*)}{C \sqrt{T}} + \frac{C (G + \epsilon)\sigma^2}{\epsilon \sqrt{T}} \nonumber \\
    &= \frac{G + \epsilon}{\sqrt{T}} \left( \frac{2(f(w_1) - f^*)}{C} + \frac{C \sigma^2}{\epsilon} \right) \tag{Step 14}
\end{align}
This proves that the average expected squared gradient norm converges to zero at a rate of $\mathcal{O}\left(\frac{1}{\sqrt{T}}\right)$.

\section*{Special Cases}
\begin{itemize}
    \item \textbf{Deterministic Setting ($\sigma^2 = 0$):} If full batch gradients are used, the variance term vanishes. In this case, we can use a constant learning rate $\alpha = \epsilon/L$, and the algorithm converges at a fast rate of $\mathcal{O}(1/T)$ to a stationary point.
    \item \textbf{No Memory ($\beta = 0$):} The algorithm simplifies to $w_{t+1} = w_t - \frac{\alpha}{|g_t| + \epsilon} \odot g_t$. This is equivalent to sign descent (or normalized gradient descent) smoothed by $\epsilon$. The proof remains mathematically identical, heavily relying on the deterministic bounds of the adaptive denominator.
    \item \textbf{Large $\epsilon$ Limit ($\epsilon \to \infty$):} As $\epsilon$ grows large, $\sqrt{v_t}$ becomes negligible, and the update rule approaches $w_{t+1} = w_t - \frac{\alpha}{\epsilon} g_t$. This recovers standard Stochastic Gradient Descent (SGD) with a fixed learning rate of $\alpha/\epsilon$.
\end{itemize}

\end{document}